% GENERATED FILE. Edit canonical lessons, then run npm run generate:books.
% canonical-source-hash: fnv1a64:2585217c555e8ec8
% canonical-lessons: FA-C250-monshi, FA-C250-mohandes, FA-C250-vakil, FA-C250-ostad, FA-C250-khalaban

\chapter{Secretary, Engineer, Lawyer, University lecturer, Pilot}
\label{ch:fa-secretary-engineer-lawyer-university-lecturer-pilot}
\hlchaptermodality{250}

\begin{chapteropening}
\textbf{By the end of this chapter:} \emph{I can say a secretary, an engineer, a lawyer, a university lecturer and a pilot.}

Complete the last lesson of chapter 250: I can say a secretary, an engineer, a lawyer, a university lecturer and a pilot.
\end{chapteropening}

\section[\texorpdfstring{\fa{منشی} (monshi)}{منشی (monshi)}]{\fa{منشی} (monshi) — a secretary}

\label{lesson:FA-C250-monshi}



\begin{quote}
Before the new one: say the Persian for celery, then the Persian for a sour cherry.
\end{quote}

\subsection*{You'll want to know: \fa{منشی}}
\textbf{\fa{منشی}} — \emph{monshi} — \textquotedblleft{}a secretary\textquotedblright{}.

\begin{cousinweb}
\emph{monshi} is borrowed from Arabic \emph{munšiʾ}, \textquotedblleft{}composer, writer\textquotedblright{}. Its Arabic root is \emph{n-š-ʾ} (Lane).
\end{cousinweb}

\begin{grammarlens}[title={What you've built}]
One more word for this chapter.
\end{grammarlens}

\subsection*{Your turn}
\begin{itemize}
\raggedright
  \item \emph{Say it:} \emph{monshi}
  \item \emph{Say it:} \emph{monshi}, once more
  \item \emph{Say it:} say \emph{monshi}
  \item \emph{Recall:} say the Persian for celery, then the Persian for a sour cherry, then say \emph{monshi} again
\end{itemize}

\begin{tcolorbox}[breakable,colback=teal!4,colframe=teal!35!black,title={Before you move on}]
What does \fa{منشی} mean? (\textquotedblleft{}A secretary\textquotedblright{}.) Say it once more.
\end{tcolorbox}
\section[\texorpdfstring{\fa{مهندس} (mohandes)}{مهندس (mohandes)}]{\fa{مهندس} (mohandes) — an engineer}

\label{lesson:FA-C250-mohandes}



\begin{quote}
Before the new one: say the Persian for a sour cherry, then the Persian for a secretary.
\end{quote}

\subsection*{You'll want to know: \fa{مهندس}}
\textbf{\fa{مهندس}} — \emph{mohandes} — \textquotedblleft{}an engineer\textquotedblright{}.

\begin{cousinweb}
\emph{mohandes} is borrowed from Arabic \emph{muhandis}, \textquotedblleft{}engineer, geometer\textquotedblright{}, built on \emph{handasa}, \textquotedblleft{}geometry\textquotedblright{}. EtymDB traces \emph{handasa} back to Persian \emph{andâze}, \textquotedblleft{}measure\textquotedblright{}: the Persian word came home in Arabic dress.
\end{cousinweb}

\begin{grammarlens}[title={What you've built}]
One more word for this chapter.
\end{grammarlens}

\subsection*{Your turn}
\begin{itemize}
\raggedright
  \item \emph{Say it:} \emph{mohandes}
  \item \emph{Say it:} \emph{mohandes}, once more
  \item \emph{Say it:} say \emph{mohandes}
  \item \emph{Recall:} say the Persian for a sour cherry, then the Persian for a secretary, then say \emph{mohandes} again
\end{itemize}

\begin{tcolorbox}[breakable,colback=teal!4,colframe=teal!35!black,title={Before you move on}]
What does \fa{مهندس} mean? (\textquotedblleft{}An engineer\textquotedblright{}.) Say it once more.
\end{tcolorbox}
\section[\texorpdfstring{\fa{وکیل} (vakil)}{وکیل (vakil)}]{\fa{وکیل} (vakil) — a lawyer}

\label{lesson:FA-C250-vakil}



\begin{quote}
Before the new one: say the Persian for a secretary, then the Persian for an engineer.
\end{quote}

\subsection*{You'll want to know: \fa{وکیل}}
\textbf{\fa{وکیل}} — \emph{vakil} — \textquotedblleft{}a lawyer\textquotedblright{}.

\begin{cousinweb}
\emph{vakil} is borrowed from Arabic \emph{wakīl}, \textquotedblleft{}agent, one entrusted\textquotedblright{}. Its Arabic root is \emph{w-k-l} (Lane).
\end{cousinweb}

\begin{grammarlens}[title={What you've built}]
One more word for this chapter.
\end{grammarlens}

\subsection*{Your turn}
\begin{itemize}
\raggedright
  \item \emph{Say it:} \emph{vakil}
  \item \emph{Say it:} \emph{vakil}, once more
  \item \emph{Say it:} say \emph{vakil}
  \item \emph{Recall:} say the Persian for a secretary, then the Persian for an engineer, then say \emph{vakil} again
\end{itemize}

\begin{tcolorbox}[breakable,colback=teal!4,colframe=teal!35!black,title={Before you move on}]
What does \fa{وکیل} mean? (\textquotedblleft{}A lawyer\textquotedblright{}.) Say it once more.
\end{tcolorbox}
\section[\texorpdfstring{\fa{استاد} (ostâd)}{استاد (ostâd)}]{\fa{استاد} (ostâd) — a university lecturer}

\label{lesson:FA-C250-ostad}



\begin{quote}
Before the new one: say the Persian for an engineer, then the Persian for a lawyer.
\end{quote}

\subsection*{You'll want to know: \fa{استاد}}
\textbf{\fa{استاد}} — \emph{ostâd} — \textquotedblleft{}a university lecturer\textquotedblright{}.

\begin{cousinweb}
\emph{ostâd} is a native Persian word, inherited rather than borrowed: it continues Middle Persian \emph{ōstād}, \textquotedblleft{}master, craftsman\textquotedblright{} (EtymDB).
\end{cousinweb}

\begin{grammarlens}[title={What you've built}]
One more word for this chapter.
\end{grammarlens}

\subsection*{Your turn}
\begin{itemize}
\raggedright
  \item \emph{Say it:} \emph{ostâd}
  \item \emph{Say it:} \emph{ostâd}, once more
  \item \emph{Say it:} say \emph{ostâd}
  \item \emph{Recall:} say the Persian for an engineer, then the Persian for a lawyer, then say \emph{ostâd} again
\end{itemize}

\begin{tcolorbox}[breakable,colback=teal!4,colframe=teal!35!black,title={Before you move on}]
What does \fa{استاد} mean? (\textquotedblleft{}A university lecturer\textquotedblright{}.) Say it once more.
\end{tcolorbox}
\section[\texorpdfstring{\fa{خلبان} (khalabân)}{خلبان (khalabân)}]{\fa{خلبان} (khalabân) — a pilot}

\label{lesson:FA-C250-khalaban}



\begin{quote}
Before the new one: say the Persian for a lawyer, then the Persian for a university lecturer.
\end{quote}

\subsection*{You'll want to know: \fa{خلبان}}
\textbf{\fa{خلبان}} — \emph{khalabân} — \textquotedblleft{}a pilot\textquotedblright{}.

\begin{cousinweb}
\emph{khalabân} is usually explained as \emph{khalâ}, \textquotedblleft{}empty space\textquotedblright{}, with \emph{-bân}, \textquotedblleft{}keeper\textquotedblright{}: the keeper of the air. \emph{khalâ} is borrowed from Arabic \emph{ḫalāʾ}, \textquotedblleft{}empty place\textquotedblright{} (root \emph{ḫ-l-w}, Lane); \emph{-bân} is native Persian.
\end{cousinweb}

\begin{grammarlens}[title={What you've built}]
One more word for this chapter.
\end{grammarlens}

\subsection*{Your turn}
\begin{itemize}
\raggedright
  \item \emph{Say it:} \emph{khalabân}
  \item \emph{Say it:} \emph{khalabân}, once more
  \item \emph{Say it:} say \emph{khalabân}
  \item \emph{Recall:} say the Persian for a lawyer, then the Persian for a university lecturer, then say \emph{khalabân} again
\end{itemize}

\begin{tcolorbox}[breakable,colback=teal!4,colframe=teal!35!black,title={Before you move on}]
What does \fa{خلبان} mean? (\textquotedblleft{}A pilot\textquotedblright{}.) Say it once more.
\end{tcolorbox}
